% !TEX TS-program = pdflatex
% !TEX encoding = UTF-8 Unicode

% This file is a template using the "beamer" package to create slides for a talk or presentation
% - Giving a talk on some subject.
% - The talk is between 15min and 45min long.
% - Style is ornate.

% MODIFIED by Jonathan Kew, 2008-07-06
% The header comments and encoding in this file were modified for inclusion with TeXworks.
% The content is otherwise unchanged from the original distributed with the beamer package.

\documentclass{beamer}
\setbeamercovered{invisible}



% Copyright 2004 by Till Tantau <tantau@users.sourceforge.net>.
%
% In principle, this file can be redistributed and/or modified under
% the terms of the GNU Public License, version 2.
%
% However, this file is supposed to be a template to be modified
% for your own needs. For this reason, if you use this file as a
% template and not specifically distribute it as part of a another
% package/program, I grant the extra permission to freely copy and
% modify this file as you see fit and even to delete this copyright
% notice. 


\mode<presentation>
{
  \usetheme{Warsaw}
  % or ...

  % or whatever (possibly just delete it)
}


\usepackage[english]{babel}
% or whatever

\usepackage[utf8]{inputenc}
\usepackage{qrcode}
\usepackage{hyperref}
\usepackage{ytableau}
% or whatever

\usepackage{times}
\usepackage[T1]{fontenc}
% Or whatever. Note that the encoding and the font should match. If T1
% does not look nice, try deleting the line with the fontenc.

%%% MATH RELATED

\title{MATH 180 Strategies of Problem Solving} 
\subtitle
{} % (optional)

\author[W.R. Casper] % (optional, use only with lots of authors)
{W.R. Casper}
% - Use the \inst{?} command only if the authors have different
%   affiliation.

\institute[California State University Fullerton] % (optional, but mostly needed)
{
  Department of Mathematics\\
  California State University Fullerton}
% - Use the \inst command only if there are several affiliations.
% - Keep it simple, no one is interested in your street address.

\subject{Talks}
% This is only inserted into the PDF information catalog. Can be left
% out. 



% If you have a file called "university-logo-filename.xxx", where xxx
% is a graphic format that can be processed by latex or pdflatex,
% resp., then you can add a logo as follows:

% \pgfdeclareimage[height=0.5cm]{university-logo}{university-logo-filename}
% \logo{\pgfuseimage{university-logo}}



% Delete this, if you do not want the table of contents to pop up at
% the beginning of each subsection:
\AtBeginSubsection[]
{
  \begin{frame}<beamer>{Outline}
    \tableofcontents[currentsection,currentsubsection]
  \end{frame}
}


% If you wish to uncover everything in a step-wise fashion, uncomment
% the following command: 

%\beamerdefaultoverlayspecification{<+->}


\begin{document}

\begin{frame}
  \titlepage
\end{frame}

\begin{frame}{Outline}
  \tableofcontents
  % You might wish to add the option [pausesections]
\end{frame}

% Since this a solution template for a generic talk, very little can
% be said about how it should be structured. However, the talk length
% of between 15min and 45min and the theme suggest that you stick to
% the following rules:  

% - Exactly two or three sections (other than the summary).
% - At *most* three subsections per section.
% - Talk about 30s to 2min per frame. So there should be between about
%   15 and 30 frames, all told.

\section{SoPS Lecture 3}
\subsection{Chomp}

\begin{frame}{Chomp}
  \begin{block}{Rules}
    \begin{itemize}
      \item create a $4\times 5$ square grid (chocolate bar)
    \end{itemize}
    \begin{center}
      \ydiagram{5,5,5,5}
    \end{center}
    \begin{itemize}
      \item turn: player takes a "bite" of the chocolate bar 
      \item bite: select a square and remove it along with any squares either above or to the right (or both)
      \item person who eats the (poisoned) last square loses
    \end{itemize}
  \end{block}
\end{frame}

\begin{frame}{Chomp}
  \begin{block}{Sample turn (P1)}
    \begin{center}
      \ydiagram{3,5,5,5}
    \end{center}
  \end{block}
\end{frame}

\begin{frame}{Chomp}
  \begin{block}{Sample turn (P2)}
    \begin{center}
      \ydiagram{3,4,4,5}
    \end{center}
  \end{block}
\end{frame}

\begin{frame}{Chomp}
  \begin{block}{Sample turn (P1)}
    \begin{center}
      \ydiagram{2,2,2,5}
    \end{center}
  \end{block}
\end{frame}

\begin{frame}{Chomp}
  \begin{block}{Sample turn (P2)}
    \begin{center}
      \ydiagram{1,1,1,5}
    \end{center}
  \end{block}
\end{frame}

\begin{frame}{Chomp}
  \begin{block}{Sample turn (P1)}
    \begin{center}
      \ydiagram{1,1,1,1}
    \end{center}
  \end{block}
\end{frame}

\begin{frame}{Chomp}
  \begin{block}{Sample turn (P2)}
    \begin{center}
      \ydiagram{0,0,0,1}
    \end{center}
  \end{block}
\end{frame}


\begin{frame}{Chomp}
  \begin{block}{Activity (20 mins)}
    \begin{itemize}
\pause
      \item split into \textbf{pairs}
\pause
      \item try plaing on a $4\times 5$ board
\pause
      \item {\color{red}record data:} did the first or second player win?
\pause
      \item {\color{red}record data:} what board states are losing positions?
\pause
      \item try to find {\color{blue}patterns}
    \end{itemize}
  \end{block}
\end{frame}

\begin{frame}{Chomp: Debrief}
\pause
  \begin{block}{Question}
    Upload a picture of the interesting losing positions that you found
  \end{block}
\pause
  \begin{block}{Submit your findings:}
    \begin{center}
      \qrcode[hyperlink,height=2cm]{https://padlet.com/williamrileycasper/pinching-pennies-losing-positions-h5xyswjj3ku3iqof}
    \end{center}
  \end{block}
\pause
  \begin{block}{Discussion:}
    Is it better to be {\color{blue}Player 1} or {\color{red}Player 2}?
  \end{block}
\end{frame}

\subsection{Strategy Stealing}

\begin{frame}{Vocabulary}
\pause
  \begin{block}{Definition}
\pause
    A \textbf{winning board} is a board state where the next player to play can force a win.\\
\pause
    A \textbf{losing board} is a board state where the next player to play can be forced to lose.
  \end{block}
\pause
  \begin{block}{Observations}
    \begin{itemize}
\pause
      \item every board is either a \textbf{winning board} or \textbf{losing board}
\pause
      \item \textbf{winning board} means there's a move to make it into a \textbf{losing board}
\pause
      \item \textbf{losing board} means \emph{every} move makes it into a \textbf{winning board}
    \end{itemize}
  \end{block}
  
\end{frame}

\begin{frame}{Examples}
  \begin{block}{Losing board}
    \begin{center}
      \ydiagram{1}
    \end{center}
  \end{block}
\end{frame}

\begin{frame}{Examples}
  \begin{block}{Winning board}
    \begin{center}
      \ydiagram{3}
    \end{center}
  \end{block}
\end{frame}

\begin{frame}{Examples}
  \begin{block}{Losing board}
    \begin{center}
      \ydiagram{1,2}
    \end{center}
  \end{block}
\end{frame}

\begin{frame}{Examples}
  \begin{block}{Winning board}
    \begin{center}
      \ydiagram{2,2}
    \end{center}
  \end{block}
\end{frame}

\begin{frame}{Player 1 Wins!}
\pause
  \begin{block}{Theorem}
    The starting rectangle is a winning board.\\
    In other words, {\color{blue} Player 1} can guarantee a win with perfect play.
  \end{block}
\pause
  \begin{block}{Why?}
    The proof strategy is called a \textbf{strategy-stealing argument}.\\
    Assume it's instead a \textbf{losing board}.
  \end{block}
\end{frame}

\begin{frame}{Strategy Stealing Argument}
\pause
  \begin{block}{Proof}
\pause
    We get to pick the first move for {\color{blue}Player 1}
\pause
    \begin{center}
      \ydiagram{4,5,5,5}
    \end{center}
\pause
    Then this should be a \textbf{winning board} for the next to play.
  \end{block}
\end{frame}

\begin{frame}{Strategy Stealing Argument}
  \begin{block}{Proof}
    There should be a move turning it into a \textbf{losing board}\\
\pause
    Guess {\color{red}Player 2}'s move:
    \begin{center}
      \ydiagram{3,3,5,5}
    \end{center}
\pause
    So if we're right, this is now a \textbf{losing board}.
  \end{block}
\end{frame}

\begin{frame}{Strategy Stealing Argument}
  \begin{block}{Proof}
    Restart the game!\\
\pause
    Have {\color{blue}Player 1} steal {\color{red}Player 2}'s strategy.\\
\pause
    Make {\color{blue}Player 1}'s first move be
    \begin{center}
      \ydiagram{3,3,5,5}
    \end{center}
\pause
    Then now is a losing board for {\color{red}Player 2}.  Impossible!
  \end{block}
\end{frame}

\begin{frame}{Strategy Stealing Argument}
  \begin{block}{Proof}
    We must have guessed wrong about {\color{red}Player 2}'s move.\\
\pause
    Re-guess {\color{red}Player 2}'s move:
    \begin{center}
      \ydiagram{3,3,3,5}
    \end{center}
\pause
    So if we're right, this is now a \textbf{losing board}.
  \end{block}
\end{frame}

\begin{frame}{Strategy Stealing Argument}
  \begin{block}{Proof}
    Restart the game!\\
\pause
    Have {\color{blue}Player 1} steal {\color{red}Player 2}'s strategy.\\
\pause
    Make {\color{blue}Player 1}'s first move be
    \begin{center}
      \ydiagram{3,3,3,5}
    \end{center}
\pause
    Then now is a losing board for {\color{red}Player 2}.  Impossible!
  \end{block}
\end{frame}

\begin{frame}{Strategy Stealing Argument}
  \begin{block}{Proof}
\pause
    Suppose {\color{red}Player 2} has a move making a \textbf{losing board}.\\
\pause
    Then {\color{blue}Player 1} could have done that as their first move instead.\\
\pause
    Therefore {\color{blue}Player 1} must be the one able to force a win!
  \end{block}
\end{frame}

\begin{frame}{Discussion}
\pause
  \begin{block}{Question}
    Does this mean that {\color{blue}Player 1}'s first move should be to take just the upper right square?
  \end{block}
\pause
  \begin{block}{Question}
    Our proof that {\color{blue}Player 1} can force a win was \textbf{non-constructive}.  We proved a winning strategy exists without finding it.  Is this satisfying?
  \end{block}
\end{frame}

\begin{frame}{Break time!}
  \begin{block}{Activity (5 minutes)}
    \begin{itemize}
      \item get up, stretch get a cup of coffee
      \item chat up your neighbor, introduce yourself
    \end{itemize}
  \end{block}
\end{frame}

\begin{frame}{Perfect Strategy}
\pause
  \begin{block}{Question}
    What is the perfect strategy for $1\times 7$ Chomp?
    \begin{center}
      \ydiagram{7}
    \end{center}
  \end{block}
\pause
  \begin{block}{Question}
    What is the perfect strategy for $1\times n$ Chomp?
    \begin{center}
      \ydiagram{7}
    \end{center}
  \end{block}
\end{frame}

\begin{frame}{Perfect Strategy}
\pause
  \begin{block}{Activity (20 mins)}
    Break up into pairs and play Chomp with different board sizes.\\
\pause
    Can you find a perfect strategy for $2\times n$ Chomp?
  \end{block}
\end{frame}

\subsection{Wrapping up}

\begin{frame}{Discussion}
\pause
  \begin{block}{Existence vs. Construction}
    In a lot of situations, we can prove something exists, long before we figure out how to find it (if ever).
    If our proof actually shows us how to find it, we call it a \textbf{constructive proof}.
  \end{block}
\pause
  \begin{block}{Question 1}
    How did our proof today compare to our previous proofs in this class?
  \end{block}
\pause
  \begin{block}{Question 2}
    Can you think of any situations where knowing something exists could be helpful, even if you don't know how to find it?
  \end{block}
\end{frame}


\begin{frame}{Question of the Day}
\pause
  \begin{block}{Puzzle}
    Start with the number $216$.\\
    On each turn, players choose a natural number which divides $216$.\\
    Once a number has been chosen, neither it nor any of it's divisors may be chosen.\\
    The player who is forced to choose $216$ loses.\\ What's the best first number to choose?
  \end{block}
\pause
  \begin{block}{Activity (10 mins)}
    Discuss the puzzle with your neighbors.  What is the solution?
  \end{block}
\end{frame}

\begin{frame}{Question of the Day}
\pause
  \begin{block}{Puzzle}
    This was just $4\times 4$ Chomp!
    $$\begin{array}{|c|c|c|c|}\hline
    8 &  4 &  2 & 1\\\hline
   24 & 12 &  6 & 3\\\hline
   72 & 36 & 18 & 9\\\hline
  216 &108 & 54 & 27\\\hline
    \end{array}$$
  \end{block}
\pause
  \begin{block}{Activity (10 mins)}
    Discuss the puzzle with your neighbors.  What is the solution?
  \end{block}
\end{frame}

\begin{frame}{Final thoughts}
  \begin{block}{Reflection}
    \begin{itemize}
      \item Look up the Intermediate Value Theorem in your calculus textbook (or online).  What is it proving exists?  Why is this proof non-constructive?
    \end{itemize}
  \end{block}
\begin{center}
Submit your responses by \textbf{Tonight!}
\end{center}
\end{frame}


\end{document}


